% GENERATED FILE. Edit canonical lessons, then run npm run generate:books.
% canonical-source-hash: fnv1a64:14b8d5b3d49e56ac
% canonical-lessons: AR-C466-ghubar, AR-C466-mughbarr, AR-C466-tasahhur, AR-C466-kathib, AR-C466-ikhtanaqa

\chapter{Dust and dunes}
\label{ch:ar-dust-and-dunes}
\hlchaptermodality{466}

\begin{chapteropening}
\textbf{By the end of this chapter:} \emph{I can say fine dust, dusty, desertification, a sand dune and to choke.}

Complete the last lesson of chapter 466: I can say fine dust, dusty, desertification, a sand dune and to choke.
\end{chapteropening}

\section[\texorpdfstring{\ar{غبار} (ghubār)}{غبار (ghubār)}]{\ar{غبار} (ghubār) — fine dust; dust in the air}

\label{lesson:AR-C466-ghubar}



\begin{quote}
Before the new one: say the Arabic for coal, then the Arabic for to consume.
\end{quote}

\subsection*{You'll want to know: \ar{غبار}}
\textbf{\ar{غبار}} — \emph{ghubār} — \textquotedblleft{}fine dust; dust in the air\textquotedblright{}.

\begin{quote}
\ar{كان} \ar{الغبار} \ar{على} \ar{كل} \ar{الكراسي}.
\end{quote}

\begin{quote}
\emph{kāna al-ghubār ʿalā kull al-karāsī.}
\end{quote}

\begin{quote}
There was dust on all the chairs.
\end{quote}

\emph{turāb} is earth and soil; \emph{ghubār} is the dust that rises and settles.

\begin{cousinweb}
\textbf{\ar{غبار}} (\emph{ghubār}) is built on the native Arabic root \emph{gh-b-r}, \textquotedblleft{}dust\textquotedblright{} (Lane).
\end{cousinweb}

\begin{grammarlens}[title={What you've built}]
One more word for getting around.
\end{grammarlens}

\subsection*{Your turn}
\begin{itemize}
\raggedright
  \item \emph{Say it:} \emph{ghubār}
  \item \emph{Say it:} \emph{ghubār}, once more
  \item \emph{Say it:} \emph{kāna al-ghubār ʿalā kull al-karāsī.} — There was dust on all the chairs.
  \item \emph{Recall:} say the Arabic for coal, then the Arabic for to consume, then say \emph{ghubār} again
\end{itemize}

\begin{tcolorbox}[breakable,colback=teal!4,colframe=teal!35!black,title={Before you move on}]
What does \ar{غبار} mean? (\textquotedblleft{}Fine dust; dust in the air\textquotedblright{}.) Say it once more.
\end{tcolorbox}
\section[\texorpdfstring{\ar{مغبر} (mughbarr)}{مغبر (mughbarr)}]{\ar{مغبر} (mughbarr) — dusty}

\label{lesson:AR-C466-mughbarr}



\begin{quote}
Before the new one: say the Arabic for to consume, then the Arabic for fine dust.
\end{quote}

\subsection*{You'll want to know: \ar{مغبر}}
\textbf{\ar{مغبر}} — \emph{mughbarr} — \textquotedblleft{}dusty\textquotedblright{}.

\begin{quote}
\ar{كان} \ar{الطريق} \ar{إلى} \ar{القرية} \ar{مغبرا}.
\end{quote}

\begin{quote}
\emph{kāna aṭ-ṭarīq ilā al-qarya mughbarran.}
\end{quote}

\begin{quote}
The road to the village was dusty.
\end{quote}

\begin{cousinweb}
\textbf{\ar{مغبر}} (\emph{mughbarr}) is built on the native Arabic root \emph{gh-b-r}, \textquotedblleft{}dust\textquotedblright{} (Lane).
\end{cousinweb}

\begin{grammarlens}[title={What you've built}]
One more word for getting around.
\end{grammarlens}

\subsection*{Your turn}
\begin{itemize}
\raggedright
  \item \emph{Say it:} \emph{mughbarr}
  \item \emph{Say it:} \emph{mughbarr}, once more
  \item \emph{Say it:} \emph{kāna aṭ-ṭarīq ilā al-qarya mughbarran.} — The road to the village was dusty.
  \item \emph{Recall:} say the Arabic for to consume, then the Arabic for fine dust, then say \emph{mughbarr} again
\end{itemize}

\begin{tcolorbox}[breakable,colback=teal!4,colframe=teal!35!black,title={Before you move on}]
What does \ar{مغبر} mean? (\textquotedblleft{}Dusty\textquotedblright{}.) Say it once more.
\end{tcolorbox}
\section[\texorpdfstring{\ar{تصحر} (taṣaḥḥur)}{تصحر (taṣaḥḥur)}]{\ar{تصحر} (taṣaḥḥur) — desertification}

\label{lesson:AR-C466-tasahhur}



\begin{quote}
Before the new one: say the Arabic for fine dust, then the Arabic for dusty.
\end{quote}

\subsection*{You'll want to know: \ar{تصحر}}
\textbf{\ar{تصحر}} — \emph{taṣaḥḥur} — \textquotedblleft{}desertification\textquotedblright{}.

\begin{quote}
\ar{التصحر} \ar{مشكلة} \ar{في} \ar{جنوب} \ar{البلد}.
\end{quote}

\begin{quote}
\emph{at-taṣaḥḥur mushkila fī janūb al-balad.}
\end{quote}

\begin{quote}
Desertification is a problem in the south of the country.
\end{quote}

\begin{cousinweb}
\textbf{\ar{تصحر}} (\emph{taṣaḥḥur}) is built on the native Arabic root \emph{ṣ-ḥ-r}, \textquotedblleft{}a desert\textquotedblright{} (Lane). It is a modern Form V verbal noun, \textquotedblleft{}becoming a desert\textquotedblright{}.
\end{cousinweb}

\begin{grammarlens}[title={What you've built}]
One more word for getting around.
\end{grammarlens}

\subsection*{Your turn}
\begin{itemize}
\raggedright
  \item \emph{Say it:} \emph{taṣaḥḥur}
  \item \emph{Say it:} \emph{taṣaḥḥur}, once more
  \item \emph{Say it:} \emph{at-taṣaḥḥur mushkila fī janūb al-balad.} — Desertification is a problem in the south of the country.
  \item \emph{Recall:} say the Arabic for fine dust, then the Arabic for dusty, then say \emph{taṣaḥḥur} again
\end{itemize}

\begin{tcolorbox}[breakable,colback=teal!4,colframe=teal!35!black,title={Before you move on}]
What does \ar{تصحر} mean? (\textquotedblleft{}Desertification\textquotedblright{}.) Say it once more.
\end{tcolorbox}
\section[\texorpdfstring{\ar{كثيب} (kathīb)}{كثيب (kathīb)}]{\ar{كثيب} (kathīb) — a sand dune}

\label{lesson:AR-C466-kathib}



\begin{quote}
Before the new one: say the Arabic for dusty, then the Arabic for desertification.
\end{quote}

\subsection*{You'll want to know: \ar{كثيب}}
\textbf{\ar{كثيب}} — \emph{kathīb} — \textquotedblleft{}a sand dune\textquotedblright{}.

\begin{quote}
\ar{صعدنا} \ar{كثيبا} \ar{كبيرا} \ar{في} \ar{الصحراء}.
\end{quote}

\begin{quote}
\emph{ṣaʿidnā kathīban kabīran fī aṣ-ṣaḥrāʾ.}
\end{quote}

\begin{quote}
We climbed a big sand dune in the desert.
\end{quote}

Plural: \ar{كثبان}, \emph{kuthbān}.

\begin{cousinweb}
\textbf{\ar{كثيب}} (\emph{kathīb}) is built on the native Arabic root \emph{k-th-b}, \textquotedblleft{}to collect together\textquotedblright{} (Lane). A dune is sand gathered into a heap.
\end{cousinweb}

\begin{grammarlens}[title={What you've built}]
One more word for getting around.
\end{grammarlens}

\subsection*{Your turn}
\begin{itemize}
\raggedright
  \item \emph{Say it:} \emph{kathīb}
  \item \emph{Say it:} \emph{kathīb}, once more
  \item \emph{Say it:} \emph{ṣaʿidnā kathīban kabīran fī aṣ-ṣaḥrāʾ.} — We climbed a big sand dune in the desert.
  \item \emph{Recall:} say the Arabic for dusty, then the Arabic for desertification, then say \emph{kathīb} again
\end{itemize}

\begin{tcolorbox}[breakable,colback=teal!4,colframe=teal!35!black,title={Before you move on}]
What does \ar{كثيب} mean? (\textquotedblleft{}A sand dune\textquotedblright{}.) Say it once more.
\end{tcolorbox}
\section[\texorpdfstring{\ar{اختنق} (ikhtanaqa)}{اختنق (ikhtanaqa)}]{\ar{اختنق} (ikhtanaqa) — to choke; to suffocate}

\label{lesson:AR-C466-ikhtanaqa}



\begin{quote}
Before the new one: say the Arabic for desertification, then the Arabic for a sand dune.
\end{quote}

\subsection*{You'll want to know: \ar{اختنق}}
\textbf{\ar{اختنق}} — \emph{ikhtanaqa} — \textquotedblleft{}to choke; to suffocate\textquotedblright{}.

\begin{quote}
\ar{اختنق} \ar{جدي} \ar{من} \ar{الغبار} \ar{في} \ar{العاصفة}.
\end{quote}

\begin{quote}
\emph{ikhtanaqa jaddī min al-ghubār fī al-ʿāṣifa.}
\end{quote}

\begin{quote}
My grandfather choked on the dust in the storm.
\end{quote}

Present: \ar{يختنق}, \emph{yakhtaniq}.

\begin{cousinweb}
\textbf{\ar{اختنق}} (\emph{ikhtanaqa}) is built on the native Arabic root \emph{kh-n-q}, \textquotedblleft{}to strangle\textquotedblright{} (Lane). It is Form VIII, \textquotedblleft{}to be choked\textquotedblright{}.
\end{cousinweb}

\begin{grammarlens}[title={What you've built}]
One more word for getting around.
\end{grammarlens}

\subsection*{Your turn}
\begin{itemize}
\raggedright
  \item \emph{Say it:} \emph{ikhtanaqa}
  \item \emph{Say it:} \emph{ikhtanaqa}, once more
  \item \emph{Say it:} \emph{ikhtanaqa jaddī min al-ghubār fī al-ʿāṣifa.} — My grandfather choked on the dust in the storm.
  \item \emph{Recall:} say the Arabic for desertification, then the Arabic for a sand dune, then say \emph{ikhtanaqa} again
\end{itemize}

\begin{tcolorbox}[breakable,colback=teal!4,colframe=teal!35!black,title={Before you move on}]
What does \ar{اختنق} mean? (\textquotedblleft{}To choke; to suffocate\textquotedblright{}.) Say it once more.
\end{tcolorbox}
